\chapter{RESULTS, DISCUSSIONS AND CONCLUSIONS} % Main chapter title
\label{ChapterResults} % For referencing the chapter elsewhere, use
% \ref{Chapter1}

\section{Results \& Analysis}
A series of case-based experiments based on Terms of Service
documents gathered from various online platforms were used to assess
the RiskWise system. Using its hybrid pipeline\index{Hybrid Retrieval}, the system processed
the uploaded documents to create knowledge graphs\index{Knowledge Graph}, produce semantic
representations, and provide risk-informed insights\index{Risk Analysis} based on the
contextual evidence gathered. Based on observations from the system
implementation and experimentation as well as the synthesis of the
literature, the results are discussed.

\subsection{Findings of Literature Survey}
The literature survey conducted as a part of this study has been very
informative in understanding the status of automated contract
analysis and research in legal AI. The following points have been
raised based on the studies surveyed:
\begin{itemize}
  \item The current state of unfair clause detection is based on
    classification models analyzing clauses independently.
  \item Most AI-based contract analysis tools lack contextual models
    that understand the relationships between clauses.
  \item Lack of explainability is a problem, and most tools are able
    to detect risks but lack the ability to explain their reasoning.
  \item The traditional Retrieval-Augmented Generation model analyzes
    text in a flat structure, which can lead to the loss of
    contextual information.
  \item Most tools are designed for legal professionals and not end-users.
\end{itemize}
Analysis:
The literature survey conducted emphasizes the need for a system that
bridges the gap between automated clause detection\index{Clause Detection} and user-friendly
explainable analysis\index{Explainable AI}. The RiskWise framework fills this gap by
incorporating hybrid retrieval based on semantic similarity and
relational reasoning.

\subsection{Results of System Implementation}

\subsection*{Contextual Retrieval Performance}
\begin{itemize}
  \item The segmentation of the uploaded documents was done into
    semantically significant chunks\index{Semantic Chunking}.
  \item Vector space representation\index{Vector Database} made it easier to retrieve
    similar clauses based on similarity.
  \item Traversal of the knowledge graph\index{Graph Traversal} helped to expose
    contractually relevant information.
  \item Hybrid retrieval combined the results to create contextual evidence\index{Contextual Retrieval}.
\end{itemize}
Result: Improved contextual completeness and efficient retrieval of
distributed contractual provisions.

\subsection*{Risk Detection and Explanation Quality}
The ToS risk identification system was tested using example ToS
documents to determine the functionality of the system in identifying
risks associated with typical ToS documents.
\begin{itemize}
  \item Unilateral modification clause identification
  \item Liability limitation clause identification
  \item Data usage and privacy permissions identification
  \item Termination and service restriction clause extraction
\end{itemize}
The explanations provided were tested for their relevance and interpretability.
Result: Risky clauses were successfully extracted along with relevant
explanations.

\subsection*{Interactive Querying Performance}
The conversational interface was tested to determine the capabilities
of user interaction and exploration of context.
\begin{itemize}
  \item Users were able to ask follow-up questions about risks identified.
  \item Contextual evidence retrieved supported informed response.
  \item Responses were still relevant to document content despite turns in
    interaction.
\end{itemize}
Result: Successful exploration of contractual risks using the
conversational interface while still grounded in context.

\subsection{Overall Observation}
The combined results demonstrate that the RiskWise system
successfully transforms complex legal agreements into structured
representations and interpretable insights. The hybrid Graph-RAG
pipeline enabled context-aware retrieval and explainable reasoning,
validating the system’s suitability for assisting users in
understanding contractual risks associated with digital services.

\section{Comparative Study}
This study compares a traditional approach for clause extraction,
text only; and a newly proposed hybrid Graph-Recurrent Attention to
Graph (RAG) Model for clause extraction. The existing literature
demonstrates that transformer based question-answering models exhibit
reasonable efficiency for clause extraction but fall short of
contextual completeness and relationship oriented reasoning.
Conversely, the proposed hybrid model in the RiskWise framework
performed better with respect to capitalising upon contextual
context, and demonstrating relationship between clauses. In addition,
the knowledge graph aspect of the model facilitated identification of
obligations/permissions which related but were not derived from
semantic similarity. Also, the hybrid model provides contextually
appropriate citation for both the text segments derived from clauses
and the evidence for the relationships, making them contextually
interpretable as well. Ultimately, this comparative analysis
demonstrates that combining structural representations of contracts
with semantic extraction improves reliability and interpretability in
automated contract analysis systems.

\section{Discussions}
The proposed project clearly demonstrates the applicability of the
knowledge graph\index{Knowledge Graph} and Retrieval-Augmented Generation\index{RAG} paradigm to be
combined for addressing legal document understanding tasks. By
combining semantic as well as relational approaches, the system
bridges the gap that has been noticed in the past. It is also
important to note that the structural representation helps in
improving the contextual completeness, especially when related
clauses are identified in different spots. Additionally, the
conversational interaction framework helps in improving accessibility
by allowing users to inquire about document understanding in natural
language queries. On the other hand, the study also points out some
intrinsic limitations of generative AI systems\index{LLM}, such as variations in
output and reliance on model capability. Other trade-offs in
performance between local and cloud inference modes have also been
noticed, which are indicative of real-world deployment aspects.

\section{Project Cost Estimation}
The RiskWise project was primarily developed using open-source
technologies and existing computational resources, resulting in
minimal direct financial expenditure. However, a conceptual cost
estimation is presented to reflect potential implementation and
operational expenses.

\begin{table}[h]
  \centering
  \caption{Project Execution Cost}
  \begin{tabular}{|c|p{2.5cm}|p{4.5cm}|c|}
    \hline
    \multicolumn{4}{|c|}{\textbf{Project Execution Cost}} \\
    \hline
    \textbf{No}  & \textbf{Category} &  \textbf{Description} &
    \textbf{Estimated Cost (Rs.)}\\
    \hline
    1  & Software Tools & Open Source Frameworks and Libraries & 0 \\
    2  & Hardware Resources &  Personal Computing resources & 0 \\
    3  & Cloud Usage &  API inference and hosting (estimated) & 2000 \\
    4  & Miscellaneous & Testing and internet usage & 1000 \\
    \hline
    \multicolumn{3}{|l|}{\textbf{Total Cost}} & 3000\\
    \hline
  \end{tabular}%
  \label{tab:executioncost}%
\end{table}%

\begin{table}[h]
  \centering
  \caption{Service Deployment Cost}
  \begin{tabular}{|c|p{4.5cm}|c|}
    \hline
    \multicolumn{3}{|c|}{\textbf{Service Deployment Cost}} \\
    \hline
    \textbf{No}  & \textbf{Component} &  \textbf{Estimated Monthly Cost(Rs.)} \\
    \hline
    1  & Cloud Hosting & 1500 \\
    2  & Model Inference API &  3000 \\
    3  & Database Hosting & 1000 \\
    \hline
    \multicolumn{2}{|l|}{\textbf{Total Cost}} & 5500\\
    \hline
  \end{tabular}%
  \label{tab:costestimation}%
\end{table}%

\section{Project Impact}

\subsection{Environmental Impact}
This application is based mainly on software and does not involve
large amounts of physical resources. Cloud-based processing may
increase indirect energy consumption. However, due to clever
retrieval mechanisms and caching strategies for the best possible
retrieval, there is less effort for computer overhead because
redundant computation is avoided.

\subsection{Societal Impact}
The significance of the RiskWise project as a whole is on the side of
humanity. The system encourages transparency and informed consent in
digital environments. The system assists users to make educated
choices regarding digital service adoption through comprehension of
terms, by reducing the complexity of contractual contracts. It helps
make a dent in consumer knowledge, enhances digital literacy, and
gives everyone equitable access to information.

\subsection{Economical Impact}
Economically, the system demonstrates potential to reduce legal
consultation costs for preliminary contract understanding and risk
identification. Organizations may also benefit from automated
contract review capabilities, improving operational efficiency. The
architecture’s reliance on open-source technologies further enhances
cost-effectiveness and accessibility.

\section{SDG Goal Compliance}
The project satisfies a number of the United Nations Sustainable
Development Goals (SDGs). The project satisfies SDG 9 (Industry,
Innovation, and Infrastructure) through the development of innovative
AI-driven infrastructure for automating document analysis. The
project also satisfies SDG 16 (Peace, Justice, and Strong
Institutions) through the promotion of transparency and fairness in
contractual relationships and through the enhancement of access to
understandable legal information. The project indirectly supports the
promotion of inclusive digital ecosystems and responsible
technological innovation through the empowerment of users with
knowledge about digital agreements.

\section{Conclusions}
The current study has been able to achieve the objectives of
designing and implementing a hybrid Graph-RAG approach for the
automated analysis of Terms of Service. The system has been able to
efficiently combine semantic embeddings\index{Semantic Embedding}, knowledge graphs\index{Knowledge Graph}, and
generative reasoning\index{Generative Reasoning} power.
The RiskWise system designed in the current research study
demonstrates the feasibility of transforming complex legal texts into
interpretable insights using context-aware analysis and reasoning.
The hybrid search strategy was able to enhance the completeness and
correctness of context awareness over the baseline approaches, and
the conversational interface was able to enhance accessibility for non-experts.
The current research study introduces a new contribution to the legal
AI community by offering a new consumer-oriented application of
hybrid search architectures.

\section{Scope for Future Work}

RiskWise sets up a base framework for automated Terms of Service analysis based on hybrid retrieval and
reasoning methodologies.
Nevertheless, there are still many research and development opportunities to improve system capability and
expand its flexibility as well.
Immediate direction is the integration of multilingual processing to support analysis of contractual documents
even in several languages.

The use of multilingual embedding models and language-agnostic large language models will ensure that the
user base will be global in nature, and also expand the usability beyond English-language agreements.
In addition, the incorpora- tion of Optical Character Recognition (OCR) processing would also allow scanned
and image-based documents to be processed and thus extending the coverage on document-based applications
and its real life relevance.

Another major extension is the ability to visualise knowledge graph structures more graphically.
Interactive and dynamic graph visualizations would enable users to explore the relationship and the implications
of contractual entity vs obligations vs permissions that makes them more interpretable and can improve
understanding in terms of how to un- derstand the document semantics.

This visualization could also be helpful to legal practitioners in exploratory contract analysis and auditing.
Personalization methods could also be researched under the model of Risk Analysis in the future.
Customization of risk profiles would enable users to set priorities for different categories such as
privacy, financial obligations, etc.

Adaptive analysis, taking into consideration the user preferences, may also prove helpful in improving the relevance of the context-aware interpretation of the risks.
The idea of browser extensions appears to be promising in the context of ToS analysis in real time.
This form of integration may prove helpful in enabling the analysis of the agreements that are encountered in the online interaction in a proactive manner before agreeing to the language of the contract.

Also, enterprise-oriented batch processing capabilities could let companies review multiple contracts with
various vendors and perform 32 comparative contract risk assessments.
Methodologically there are opportunities to broaden the reasoning ability of the enterprise, by including legal
terminol- ogy and databases for their domain knowledge.

For example, the inclusion of external legal standards and legal precedent data in the augmentation of the building of the knowledge graph could potentially allow for the development of semantic richness which could allow for the development of complex reasoning in compliance.
The focus of the research agenda will also include the participation in the empirical validation of the proposed framework through the conduct of quantitative benchmarking against annotated legal datasets as well as user performance in relation to metrics of system performance in relation to the identification of clauses that can be achieved in relation to the retrieval of the same information.
User studies can also be included in the identification of user performance in relation to usability, trust and interpretability.

The potential to explore and research possibilities for distributed infrastructures in data processing, optimally indexed storage technologies and compression technologies would help to assist application opportunities.
The potential to improve privacy and security through device based inferences and encrypted knowledge storage would improve user experiences.
RiskWise is a scalable platform and will continue to enable research into AI-enhanced contract comprehension.

